\chapter*{Bibliographie}
\addcontentsline{toc}{chapter}{Bibliographie}

\begin{description}[style=nextline,leftmargin=2.2cm,labelwidth=2cm]
  \item[{[BF]}] R. L. Burden, J. D. Faires, A. M. Burden, \textit{Numerical
    Analysis}, 10\textsuperscript{e} édition, Cengage Learning, 2016. Les
    renvois donnent les numéros d'algorithmes, de théorèmes ou d'équations et
    les pages imprimées du livre.
  \item[{[Cours]}] P. Woafo, \textit{Méthodes numériques pour physiciens et
    ingénieurs}, Université de Yaoundé I (transcription LaTeX dans le dossier
    \fichier{latex/} de ce dépôt).
  \item[{[OEIS]}] The On-Line Encyclopedia of Integer Sequences, suite A094090
    (solution positive de $5(1 - e^{u}) + ue^{u} = 0$, constante de la loi de
    Wien), \url{https://oeis.org/A094090}.
  \item[{[VdW]}] Purdue University, \og Deviations from the Ideal Gas Law \fg{}
    (table des constantes de Van der Waals),
    \url{https://chemed.chem.purdue.edu/genchem/topicreview/bp/ch4/deviation.php}.
  \item[{[CSD]}] Équation de Carnahan-Starling pour les sphères dures :
    \og Playing with Marbles: Structural and Thermodynamic Properties of
    Hard-Sphere Systems \fg, \url{https://arxiv.org/abs/1310.5578} ;
    équation de Carnahan-Starling-De Santis :
    \url{https://iifiir.org/en/fridoc/a-modified-carnahan-starling-de-santis-equation-of-state-to-compute-106555}.
  \item[{[CL]}] Loi de Child-Langmuir : \og Child's Law \fg, SIMION,
    \url{https://simion.com/info/childs_law.html} ; \og A new approach to the
    Child-Langmuir law \fg, \url{https://arxiv.org/abs/1506.07417}.
  \item[{[LA]}] F. M. S. Lima, P. Arun, \og An accurate formula for the period of
    a simple pendulum oscillating beyond the small-angle regime \fg,
    \url{https://arxiv.org/abs/physics/0510206} (formule
    $T = \frac{2T_0}{\pi}K(\sin\frac{\theta_0}{2})$).
  \item[{[TW]}] L. N. Trefethen, J. A. C. Weideman, \og The exponentially
    convergent trapezoidal rule \fg, \textit{SIAM Review} 56(3), 2014,
    p.~385--458.
  \item[{[VN]}] \og Von Neumann stability analysis \fg, Wikipédia (anglais),
    \url{https://en.wikipedia.org/wiki/Von_Neumann_stability_analysis}
    (stabilité inconditionnelle du schéma de Dufort-Frankel).
  \item[{[DF2]}] \og New analysis of the Du Fort-Frankel methods \fg,
    \textit{Journal of Scientific Computing}, 2012,
    doi:10.1007/s10915-012-9627-2,
    \url{https://cris.iucc.ac.il/en/publications/new-analysis-of-the-du-fort-frankel-methods/}
    (erreur de troncature $O(k^2 + h^2 + k^2/h^2)$).
  \item[{[HGG]}] A. C. Hindmarsh, P. M. Gresho, D. F. Griffiths, \og The
    stability of explicit Euler time-integration for certain finite difference
    approximations of the multi-dimensional advection-diffusion equation \fg,
    \textit{International Journal for Numerical Methods in Fluids} 4(9), 1984,
    doi:10.1002/fld.1650040905.
\end{description}

\section*{Où trouver d'autres exercices corrigés}

Les exercices du cours reprennent des exercices classiques. Le livre [BF]
donne les réponses d'une sélection d'exercices à la fin de l'ouvrage
(\og Answers for Selected Exercises \fg, p.~787 et suivantes). En français,
l'ouvrage d'A. Quarteroni, F. Saleri et P. Gervasio, \textit{Calcul
scientifique : cours, exercices corrigés et illustrations en MATLAB et Octave}
(Springer), traite les mêmes méthodes avec des exercices corrigés.
